\title{First Steps Toward Automated AI Research - Recursive}
\begin{document}
\maketitle

\begin{abstract}
22.06.26, 15:53 First Steps Toward Automated AI Research - Recursive Seite 1 von 14 https://www.recursive.com/articles/first-steps-toward-automated-ai-research?utm\_source=substack\&utm\_medium=email First Steps Toward Automated AI Research Early results from Recursive’s automated AI research system on model training and GPU kernel benchmarks JUNE 11, 2026 Today we are releasing early results from Recursive’s automated AI research system. Across three benchmarks, the system achieves state-of-the-a! results: in "xed-budget language model training, small-model training speed, and GPU kernel optimization. The system automates the research loop for a target objective: it proposes an idea, implements it, runs an experiment, validates the result, and uses what it learns to choose the next experiment. It runs many research threads over long horizons, keeps useful context from prior experiments, combines promising branches, and puts results through validation for reward hacks and variance 22.06.26, 15:53 First Steps Toward Automated AI Research - Recursive Seite 2 von 14 https://www.recursive.com/articles/first-steps-toward-automated-ai-research?utm\_source=substack\&utm\_medium=email before treating improved pe\#ormance as real progress. It is designed to scale and harnesses principles of open-ended algorithms, building on ideas from previous work by our team and others into recursively self-improving AI. We tested the system on benchmarks chosen for both practical impo!ance and tight feed
\end{abstract}

\section{Recovered Text}
22.06.26, 15:53
First Steps Toward Automated AI Research - Recursive
Seite 1 von 14
https://www.recursive.com/articles/first-steps-toward-automated-ai-research?utm\_source=substack\&utm\_medium=email
First Steps Toward
Automated AI Research
Early results from Recursive’s automated AI research system
on model training and GPU kernel benchmarks
JUNE 11, 2026
Today we are releasing early results from Recursive’s automated AI research
system. Across three benchmarks, the system achieves state-of-the-a!
results: in "xed-budget language model training, small-model training
speed, and GPU kernel optimization.
The system automates the research loop for a target objective: it proposes an
idea, implements it, runs an experiment, validates the result, and uses what it
learns to choose the next experiment. It runs many research threads over long
horizons, keeps useful context from prior experiments, combines promising
branches, and puts results through validation for reward hacks and variance
22.06.26, 15:53
First Steps Toward Automated AI Research - Recursive
Seite 2 von 14
https://www.recursive.com/articles/first-steps-toward-automated-ai-research?utm\_source=substack\&utm\_medium=email
before treating improved pe\#ormance as real progress. It is designed to
scale and harnesses principles of open-ended algorithms, building on ideas
from previous work by our team and others into recursively self-improving AI.
We tested the system on benchmarks chosen for both practical impo!ance and
tight feedback loops. They stress three core levers of AI progress: be\$er training
algorithms, faster training, and more e\%cient use of hardware. They are also
well suited to automated research because they have clear metrics, relatively
low variance, and evaluators that can be hardened against reward hacks.
We are open-sourcing a!ifacts from these runs so others can inspect
and build on the system’s outputs.
Benchmark
Task Type
Metric
Previous State
of the A"
Recursive
Improvem
NanoChat
Autoresearch
Train a small
language
model to
highest
pe!ormance
given a small
compute
budget
Validation
BPB
0.9372
0.9109
0.0263 lo
Validation
BPB, or a
speedup
reach the
same loss
NanoGPT
Speedrun
Train a small
language
model to a
ce"ain
pe!ormance
as fast as
possible
Training
time
required
to reach
a 3.28
validation
loss
79.7 s
77.5 s
2.2s faste
training
SOL-
ExecBench
Optimize
GPU kernels
toward
Mean
SOL
score
0.699
0.754
18\%
reduction
gap to th
22.06.26, 15:53
First Steps Toward Automated AI Research - Recursive
Seite 3 von 14
https://www.recursive.com/articles/first-steps-toward-automated-ai-research?utm\_source=substack\&utm\_medium=email
Case study 1: NanoChat Autoresearch
Andrej Karpathy’s NanoChat autoresearch repo is a popular sta!ing point for
automated research systems. The task is to train a small language model to the
lowest validation loss, measured in bits per byte (BPB), within a "xed "ve-minute
budget on a single GPU. It is a natural test of our system because experiments
are fast, variance is low, and reward hacks are relatively easy to detect.
Perhaps for those reasons, a public collaborative e\&o! has already formed
around this setup. autoresearch@home extends the original setup into a
collaborative se\$ing where several dozens of humans and hundreds of their
agents collectively improve pe\#ormance. That gives us a stronger comparison
point than Karpathy’s single overnight run. We wanted to test if our system could
improve on solutions produced by an entire community of humans and agents.
Our system sta!s from the same initial seed solution the Autoresearch code
sta!s from. We initially searched on NVIDIA H100 GPUs, then transferred the
discovered solution to run on an NVIDIA B200 GPU for a fair comparison to
public results. After removing minor reward hacks from the previous best
autoresearch@home solution and evaluating it on 10 random seeds, its
mean pe\#ormance is 0.9372 BPB. Our system found a solution that reached
0.9109 BPB, a 0.0263 BPB improvement. Measured another way, our solution
reaches the quality of Karpathy’s original overnight autoresearch BPB in
roughly 1.3x less training time than the best autoresearch@home solution.
hardware
limits
across
235
kernels
optimal
pe!orma
estimate
1.0
22.06.26, 15:53
First Steps Toward Automated AI Research - Recursive
Seite 4 von 14
https://www.recursive.com/articles/first-steps-toward-automated-ai-research?utm\_source=substack\&utm\_medium=email
FIGURE 1
Autoresearch sta!s from an already optimized model with some non-trivial
design decisions baked in. To this end, we tested whether our system could
also make improvements from a much weaker sta!ing point, a naive initial
implementation (a vanilla Transformer with AdamW). Our system improved
the model from 1.059 BPB to 0.9344 BPB (evaluated on an NVIDIA B200 GPU),
again outpe\#orming the best solution produced by the autoresearch@home
community. This does not necessarily prove independent rediscovery, since
the underlying models may know many public techniques including those
used by or created by the autoresearch@home community, but it does
show that the search process can assemble a competitive training stack
from a much weaker sta!ing point. The resulting solution also di\&ered
in several ways from the public best solution.
22.06.26, 15:53
First Steps Toward Automated AI Research - Recursive
Seite 5 von 14
https://www.recursive.com/articles/first-steps-toward-automated-ai-research?utm\_source=substack\&utm\_medium=email
FIGURE 2
FIGURE 3
What modi"cations did our system come up with? The best solutions were
22.06.26, 15:53
First Steps Toward Automated AI Research - Recursive
Seite 6 von 14
https://www.recursive.com/articles/first-steps-toward-automated-ai-research?utm\_source=substack\&utm\_medium=email
not driven by one trick. They combined changes to architecture, sho!-
context memory, auxiliary losses, a\$ention, optimizer behavior, weight decay
schedules, compiler se\$ings, and more.
One of the biggest gains came from a richer sho!-context memory mechanism.
The baseline already uses value embeddings; our system extended this
idea with hashed bigram and trigram embedding tables, mixed into the
a\$ention value path through learned gates. This gave the model a cheap
way to use local n-gram information without paying the time cost of slower
convolutional or a\$ention-heavy alternatives.
This connects to recent work such as DeepSeek Engram, which explores
hash tables as a sparsity axis. In our se\$ing, hash tables can add 1-2 billion
sparse parameters to a roughly 50M parameter model: most entries are
inactive on any given batch, and lookup is cheap. Similar hash-table and
n-gram ideas also appear in top NanoGPT Speedrun submissions. The
system adapted this family of ideas to the "xed-budget se\$ing by injecting
hashed bigram and trigram embeddings into a\$ention value vectors across
multiple layers, with di\&erent hashes per layer to reduce repeated collisions.
We are not aware of prior work using this exact variant.
The expandable boxes below include selected technical details from the
system’s solutions. We manually inspected these outputs and used AI-
assisted analysis to understand the techniques and screen for reward hacks.
We may still have missed errors in kernel optimization where we are not
specialists, but that is also pa! of the point: the ideas presented here came
from the system, not from our prior expe!ise.
Hash tables
22.06.26, 15:53
First Steps Toward Automated AI Research - Recursive
Seite 7 von 14
https://www.recursive.com/articles/first-steps-toward-automated-ai-research?utm\_source=substack\&utm\_medium=email
The run optimizing the vanilla Transformer used some of the same techniques
as our best solution, including hash tables and squared-ReLU MLPs. But it also
converged on a di\&erent (yet equally competitive) "nal stack, including token-
shifting, weight averaging before eval, and byte-level feature embeddings.
This suggests the system was not merely repeating the same discoveries it
found in the other run. The expandable box below shows a few modi"cations
unique to the vanilla Transformer run.
NanoChat shows how asking our system to improve "xed-budget training led to
the discovery of many compounding, budget-aware improvements. The next test
was whether the same process could still "nd gains after years of public human
optimization on a benchmark. We tested that on NanoGPT Speedrun, whose
best public solution has been highly optimized by the community over two years.
Case study 2: NanoGPT Speedrun
NanoGPT Speedrun is a similar task, yet it's much harder to beat the state
of the a! because a large community has been optimizing solutions for
it for over two years. Instead of asking how low of a validation loss can
be achieved in a "xed time budget, the benchmark asks how quickly a
small GPT-style model can be trained to a "xed validation loss of 3.28 on
the FineWeb text dataset, using a single HGX H100 8-GPU node.
This is a mature community e\&o!, with 83 human record-se\$ing contributions
to the leaderboard so far and hundreds of proposed PRs. Since mid-2024,
the training time has been pushed from roughly 45 minutes down to 79.7
Optimizing a vanilla Transformer
22.06.26, 15:53
First Steps Toward Automated AI Research - Recursive
Seite 8 von 14
https://www.recursive.com/articles/first-steps-toward-automated-ai-research?utm\_source=substack\&utm\_medium=email
seconds through a long sequence of primarily hand-engineered submissions.
Given that the current solution is so well optimized, there are few obvious
improvements left.
Sta!ing from the current leading solution, our system found a set of additional
optimizations that reduced training time from 79.7 seconds to 77.5 seconds
while still meeting the leaderboard’s validation-loss signi"cance requirement
(mean validation loss ≤ 3.28 at p < 0.01).  This is a similar or larger improvement
than recent human contributions.
FIGURE 4
We also tested whether the system could make progress from a much weaker
sta!ing point. Sta!ing from an earlier roughly 15-minute solution (we chose the
1
The 77.5 s solution: FP8 a!ention,
exploration noise, and cautious embeddings
22.06.26, 15:53
First Steps Toward Automated AI Research - Recursive
Seite 9 von 14
https://www.recursive.com/articles/first-steps-toward-automated-ai-research?utm\_source=substack\&utm\_medium=email
"rst working Python solution in the lineage of champions (entry \#5), as earlier
versions were in C), our system reached approximately 185 seconds in a few
days, close to the human leaderboard’s May 2025 roughly 180-second level. This
should not be read as independent or unique discovery, since the underlying
models may have seen the repository, but the system found a di\&erent
"nal solution and added the overlapping contributions in a di\&erent order.
The 77.5s solution was not a single optimization. It combined changes to
a\$ention precision, optimizer behavior, embedding updates, schedule
choices, and fused GPU kernels. Each change had to save time without
destabilizing training.
Despite an entire community of humans (sometimes with AI assistance)
spending years working on this problem, Recursive’s automated AI research
system still discovered additional improvements. The next case study moves
one level lower, from small-model training recipes to GPU kernels. Unlike the
"rst two benchmarks, kernel optimization is closer to production systems
work: it often determines the cost of real training and inference workloads.
Case study 3: SOL-ExecBench
The "rst two benchmarks optimize small language model training runs. SOL-
ExecBench instead focuses on writing fast, correct GPU kernels: the small
accelerator programs behind operations such as matrix multiplications,
reductions, normalization layers, a\$ention components, quantization routines,
and fused blocks.
Analysis of the 185 s from-scratch
solution: a di"erent route to 3 minutes
22.06.26, 15:53
First Steps Toward Automated AI Research - Recursive
Seite 10 von 14
https://www.recursive.com/articles/first-steps-toward-automated-ai-research?utm\_source=substack\&utm\_medium=email
The benchmark contains 235 kernel-writing tasks derived from real workloads.
Each task provides a simple reference PyTorch implementation that de"nes
the signature, tensor shapes, data types, and numerical contract (what
output the kernel must produce, and how close it must be to the reference
implementation). The goal is to produce the same result within tolerance
while running as fast as possible on NVIDIA Blackwell B200 GPUs.
The benchmark repo!s a Speed-of-Light (SOL) SOL-ExecBench score: 0.5
corresponds to the benchmark’s optimized PyTorch baseline, and a score of
1.0 corresponds to the benchmark’s analytical optimal pe\#ormance estimate.
We ran our system across all 235 kernels jointly, so it could reuse its discoveries
for be\$er ways to do things across related tasks (e.g. pa\$erns for memory
movement, tiling, reductions, vectorization, and fusion). We provided standard
pro"ling tools but did not speci"cally tune the system for kernel engineering.
Aside from adding pro"ling tools, we use the same system to optimize
kernels as in the other two benchmarks.
Our system achieved a mean NVIDIA SOL-ExecBench score of 0.754, an 18\%
reduction in the gap to the hardware limit from the previous leaderboard
best of 0.699.
22.06.26, 15:53
First Steps Toward Automated AI Research - Recursive
Seite 11 von 14
https://www.recursive.com/articles/first-steps-toward-automated-ai-research?utm\_source=substack\&utm\_medium=email
FIGURE 5
FIGURE 6
We checked a few high-pe\#orming kernels and found that the solutions include
1.0
0.9
SOLscore
SOL-ExecBench:SOLscoreperkernel
Recursivevs.leaderboardbest•higherisbetter•0.5=optimizedPyTorch,1.0=hardwarerooflin
0.8
0.7
MeanSOLscore
SOL-ExecBench:meanSOLscorebykernelcategory
Recursivevs.leaderboardbest,doubleAl,andCursor•higherisbetter(zoomedto0.5-1.0)
1.0
•Recursive
Leaderboardbest
doubleAl
0.9
0.809
0.8
0.761
-0.7460.731
0,758
0.724
0.7160.709
0.7
0 637
0.77
22.06.26, 15:53
First Steps Toward Automated AI Research - Recursive
Seite 12 von 14
https://www.recursive.com/articles/first-steps-toward-automated-ai-research?utm\_source=substack\&utm\_medium=email
a range of good kernel engineering practices and creative solutions. The
expandable boxes provide a few examples for those interested in the details.
While reward hacking was an issue we contended with for all three benchmarks,
it was pa!icularly challenging on SOL-ExecBench. Some candidates exploited
the evaluation setup instead of implementing genuinely faster kernels:
caching outputs, relying on persistent state, or taking advantage of
timing-harness details.
For that reason, we treated a correctness audit as pa! of the research
system on all of the benchmarks. Promising improvements were passed
through increasingly strict automated checks designed to distinguish genuine
kernel improvements from benchmark-speci"c exploits. This substantially
reduced reward hacks and became an impo!ant pa! of the loop itself: as
the search became stronger, the evaluator had to become stronger too.
SOL-ExecBench demonstrates the ability of our system to improve an entirely
Example L1 kernel
045\_fused\_linear\_gelu\_grn\_linear (L1)
Example L2 kernel
049\_group\_limited\_topk\_routing (L2)
Example Quant kernel
032\_nvfp4\_moe\_expe!\_linear\_with\_gating (Quant)
Example FlashInfer-Bench kernel
012\_gqa\_paged\_decode\_h32\_kv4\_d128\_ps1 (FlashInfer-Bench)
22.06.26, 15:53
First Steps Toward Automated AI Research - Recursive
Seite 13 von 14
https://www.recursive.com/articles/first-steps-toward-automated-ai-research?utm\_source=substack\&utm\_medium=email
di\&erent pa! of the AI stack. It had to reason about low-level implementation
choices, generate candidate kernels, run correctness and pe\#ormance
checks, and transfer useful pa\$erns across related tasks.
What’s Next
These results are an early sign that our system can push the frontier on AI
training and infrastructure tasks, especially when the goal is well-de"ned,
measurable, and quick enough to evaluate many times. The system made
progress by compounding many discoveries: inventing new optimizations,
recasting known ideas under tighter constraints, tuning implementation
details that ma\$ered, and composing improvements across modeling,
optimization, and systems layers.
Throughout this work, and especially as our search becomes more powe\#ul,
a key challenge is reward hacking (i.e. making sure the system solves the
intended task instead of exploiting loopholes that meet the le\$er of the task
and score highly, but subve! the intention of the task). We implemented many
techniques to avoid and detect such reward hacking, including iteratively
improving a reward hacking detector with AI-assisted and/or human feedback.
We expect this will remain necessary as we tackle ever more challenging
real-world applications and create more powe\#ul automated AI research
algorithms. Aligning such systems to solve the spirit of the task, and not its
le\$er, will be a grand challenge of creating systems that automate knowledge
discovery and recursively self-improve in a way that is safe and helpful.
We are excited to continue to work on that essential problem.
Many of the gains here improve e\%ciency. That ma\$ers because AI progress
does not come only from larger models and more compute; it also comes
22.06.26, 15:53
First Steps Toward Automated AI Research - Recursive
Seite 14 von 14
https://www.recursive.com/articles/first-steps-toward-automated-ai-research?utm\_source=substack\&utm\_medium=email
from making existing systems train faster, run cheaper, and use hardware
more e\&ectively. We expect systems like this to reduce the cost of intelligence:
"rst by "nding be\$er engineering tradeo\&s in today’s systems, and over
time by automating larger pa!s of the frontier research process itself.
We are open-sourcing a!ifacts from these runs so others can inspect
and build on the system’s outputs. If you’re interested in building systems
that make automated research more capable and bene"cial for humanity,
please apply to join us.
FOOTNOTES
1. We obtained our results on Modal HGX H100 8-GPU nodes and independently re-
con"rmed the numbers on Andromeda HGX H100 8-GPU nodes within noise. We
are awaiting access to PrimeIntellect HGX H100 8-GPU nodes (the o\#cial hardware)
to submit to the leaderboard.
© 2026 Recursive Superintelligence, Inc.
All rights reserved.
Terms
of use
Privacy PolicyFollow
us
LinkedIn
X
\end{document}
